\section{Resultados}
    En esta sección se presentan los resultados y mejoras visuales aplicadas 
    en cada imagen del dataset propuesto.

    \subsection{Imagen 1}
        En esta primera imagen se puede evidenciar que tiene un bajo brillo, por lo que es 
        conveniente la aplicación de la corrección gamma o de la transformación gain-bias 
        para incrementar su iluminación general. De la escala de grises, se puede apreciar 
        que su contraste es el mínimo para distinguir sus elementos más no el más adecuado, 
        por lo que también se podría mejorar. 

        \begin{figure}[H]
            \centering
            \includegraphics[width=0.9\textwidth]{../results/describe_imagen1.jpg}
            \caption{Descripción visual de la Imagen 1 en la que se puede observar un problema de brillo y de contraste a mejorar.}
        \end{figure}

        \begin{figure}[H]
            \centering
            \includegraphics[width=\textwidth]{../results/channels_imagen1.jpg}
            \caption{Histogramas de canales de la Imagen 1 donde se evidencia el bajo brillo de la imagen y la posibilidad de mejorar el contraste.}
        \end{figure}

        Al aplicar la corrección gamma con un gamma menor a 1 propicia que la imagen tienda 
        a colores oscuros debido a que los propios valores se ven reducidos. Por ello, al 
        emplear un gamma mayor a 1 hace que la imagen se vea más ``viva'' y, lo más importante, 
        se distingan los colores adecuadamente. De igual manera, usar un gamma muy alto (mayor a 
        1.4) hace que todos las intensidades tiendan a valores cada vez más próximos a 255, 
        convirtiendo totalmente en blanco a la imagen.

        \begin{figure}[H]
            \centering
            \includegraphics[width=\textwidth]{../results/gammas_imagen1.jpg}
            \caption{Corrección gamma a la Imagen 1 en la que un valor de gamma mayor a 1 mejora la calidad de la imagen.}
        \end{figure}

        Empleando la transformación gain-bias se logra los mismos resultados que por la 
        corrección gamma, es decir, una mejora en el brillo final de la imagen. Por lo que 
        ambas técnicas logran los resultados esperados sin tener que realizar una búsqueda 
        exhaustiva de los valores en las transformaciones.

        % \begin{figure}[H]
            % \centering
            % \includegraphics[width=\textwidth]{../results/gain_bias_imagen1.jpg}
            % \caption{Transformación gain-bias de la Imagen 1 con una ganancia positiva mejora la iluminación y distinción de los colores en la imagen resultante.}
        % \end{figure}

    \subsection{Imagen 2}
        En esta segunda imagen se puede mostrar que tiene un alto brillo, por lo que una transformación 
        gain-bias podría mejorar la calidad de la imagen sin alterar su contraste o colores. De la 
        escala de grises, se pueden delimitar, en la medida de lo posible, los elementos que componen a 
        imagen, esto sugiere que el problema subyacente es el brillo mismo.\\

        \begin{figure}[h]
            \centering
            \includegraphics[width=0.9\textwidth]{../results/describe_imagen2.jpg}
            \caption{Descripción visual de la Imagen 2 en la que se puede observar un problema de mucho brillo.}
        \end{figure}       

        \begin{figure}[h]
            \centering
            \includegraphics[width=\textwidth]{../results/channels_imagen2.jpg}
            \caption{Histogramas de canales de la Imagen 2 donde se evidencia el alto brillo de la imagen y un contraste deficiente.}
        \end{figure}

        Al aplicar la corrección gamma, con valores inclusive ligeramente diferentes a 1, se puede 
        apreciar que la calidad de la imagen empeora de forma drástica, este hecho muestra que 
        esta técnica no es la más adecuada para reducir el impacto del brillo sin deformar las 
        tonalidades y el propio contraste de la imagen. \\

        % \begin{figure}[H]
            % \centering
            % \includegraphics[width=\textwidth]{../results/gammas_imagen2.jpg}
            % \caption{Corrección gamma a la Imagen 2 en la que un valor de gamma menor a 1 podría mejorar la calidad, mas no de la forma esperada.}
        % \end{figure}

        Como se mencionó, la transformación gain-bias logra los resultados y mejoras esperadas, en el 
        sentido de que se reduce el brillo lo suficiente para no reducir tampoco el contraste. Lo cual 
        favorece a que se puedan distinguir de mejor manera los elementos que componen a la imagen, 
        por lo que resulta más conveniente esta técnica que la corrección gamma. Para lograr este 
        efecto fue necesario realizar un estimación de los valores más exhaustiva debido a que usar 
        un valor de ganancia mayor a 1.4 provocaba que la imagen se volviese blanca, y también se tuvo 
        que ajustar esa sobreexposición un sesgo conveniente. 

        \begin{figure}[H]
            \centering
            \includegraphics[width=\textwidth]{../results/gain_bias_imagen2.jpg}
            \caption{Transformación gain-bias de la Imagen 2 con un sesgo negativo hace posible bajar la iluminación de la imagen y mejorar ligeramente el contraste de la imagen final.}
        \end{figure}

    \subsection{Imagen 3}

        En la tercera imagen se muestra una clara ausencia de iluminación, la cual impede 
        vislumbrar los elementos presentes. Por la escala de grises y los mismos histogramas 
        sugieren que se ocupa de incrementar tanto el brillo como el contraste para lograr 
        una mejora en la calidad de la imagen. Por lo que la corrección gamma va a permitir 
        este incremento de forma natural sin distorsionar los colores.

        \begin{figure}[H]
            \centering
            \includegraphics[width=0.9\textwidth]{../results/describe_imagen3.jpg}
            \caption{Descripción visual de la Imagen 3 en la que se puede observar un problema de brillo a mejorar.}
        \end{figure}

        \begin{figure}[H]
            \centering
            \includegraphics[width=\textwidth]{../results/channels_imagen3.jpg}
            \caption{Histogramas de canales de la Imagen 3 donde se evidencia el bajo brillo de la imagen y un contraste deficiente.}
        \end{figure}

        Como se mencionó, al aplicar una corrección gamma con un gamma mayor a 1 se logra iluminar 
        la imagen sin sobreexponerla mientras se preserva un contraste más natural. En el caso contrario, 
        con una gamma menor a 1, la imagen se pierde en colores más oscuros, y con un mayor a 1.35 
        empiezan a aparecer tonalidades más naranjadas o zonas con iluminación no natural.

        \begin{figure}[H]
            \centering
            \includegraphics[width=\textwidth]{../results/gammas_imagen3.jpg}
            \caption{Corrección gamma a la Imagen 3 en la que un valor de gamma mayor a 1 podría mejorar la calidad de la forma esperada.}
        \end{figure}

        Así como ocurrió en la imagen 1, con adecuados valores de ganancia y sesgo se pueden 
        lograr resultados parecidos a los generados por una corrección gamma, pero esta técnica 
        es mas eficientes para estos casos que se requieren incrementar tanto el brillo como el 
        contraste. 

        % \begin{figure}[H]
            % \centering
            % \includegraphics[width=\textwidth]{../results/gain_bias_imagen3.jpg}
            % \caption{Transformación gain-bias de la Imagen 2 con un sesgo negativo hace posible bajar la iluminación de la imagen y mejorar ligeramente el contraste de la imagen final.}
        % \end{figure}

    \subsection{Imagen 4}

        En esta cuarta imagen se caracteriza por tener una zona central sobreexpuesta pero 
        cuenta con un buen contraste. Es decir, por su escala de grises se puede observar 
        los diferentes elementos de la imagen y se reduce a que su principal problema es 
        un brillo alto, que se podría resolver con una adecuada transformación gain-bias.\\

        \begin{figure}[h]
            \centering
            \includegraphics[width=0.9\textwidth]{../results/describe_imagen4.jpg}
            \caption{Descripción visual de la Imagen 4 en la que se puede observar un problema de iluminación en la zona central.}
        \end{figure}

        \begin{figure}[h]
            \centering
            \includegraphics[width=\textwidth]{../results/channels_imagen4.jpg}
            \caption{Histogramas de canales de la Imagen 4 donde se muestra que cuenta con un buen contraste.}
        \end{figure}

        Al aplicar una corrección gamma con valores ligeramente diferentes a 1 hacen que se 
        degrade la calidad de la imagen en el sentido de que se oscurece o ilumina 
        demasiado, esto por tener histogramas bien ecualizados y mueva demasiados valores 
        a los extremos de intensidad.\\

        % \begin{figure}[H]
            % \centering
            % \includegraphics[width=\textwidth]{../results/gammas_imagen4.jpg}
            % \caption{Corrección gamma a la Imagen 2 en la que un valor de gamma menor a 1 podría mejorar la calidad, mas no de la forma esperada.}
        % \end{figure}

        En cambio, al emplear la transformación gain-bias se logra una reducción del brillo 
        de la zona central sin oscurece toda la imagen de forma natural. Esto se debe a que 
        esta transformación permite realizar cambios más suaves sin tener que alterar 
        drásticamente todos los colores. 

        \begin{figure}[H]
            \centering
            \includegraphics[width=\textwidth]{../results/gain_bias_imagen4.jpg}
            \caption{Transformación gain-bias de la Imagen 4 con una ganancia positiva y menor a 1 permite reducir el brillo de la zona central.}
        \end{figure}

    \subsection{Imagen 5}

        La quinta imagen presenta una evidente falta en la intensidad del canal verde, lo cual afecta 
        directamente en como es percibida la imagen, viéndose extraña o poco natural. Por lo tanto, la 
        transformación se debería que aplicar únicamente al canal verde, específicamente ambas técnicas 
        van a permitir mejorar la presencia de este canal.
        
        \begin{figure}[H]
            \centering
            \includegraphics[width=0.9\textwidth]{../results/describe_imagen5.jpg}
            \caption{Descripción visual de la Imagen 5 en la que se puede observar un problema de subeexposición en el canal verde.}
        \end{figure}

        \begin{figure}[H]
            \centering
            \includegraphics[width=\textwidth]{../results/channels_imagen5.jpg}
            \caption{Histogramas de canales de la Imagen 5 donde se evidencia que tiene un buen contraste más no una buena representación del canal verde.}
        \end{figure}

        Si se aplica una corrección gamma sobre el canal verde con un gamma menor a 1, provocará 
        una reducción del color verde en la imagen. Mientras con un valor mayor a 1, se logra el 
        efecto deseado de incrementar su intensidad, así realizando una corrección significativa.

        \begin{figure}[H]
            \centering
            \includegraphics[width=\textwidth]{../results/gammas_imagen5.jpg}
            \caption{Corrección gamma a la Imagen 5 en la que un valor de gamma mayor a 1 sobre el canal verde logra la mejora deseada.}
        \end{figure}

        De igual forma, la transformación gain-bias sobre el canal verde permite la corrección de 
        este color logrando una mejora en la calidad visual resultante. Por lo que ambas técnicas 
        logran los efectos positivos sobre la imagen. 

        % \begin{figure}[H]
            % \centering
            % \includegraphics[width=\textwidth]{../results/gain_bias_imagen5.jpg}
            % \caption{Transformación gain-bias de la Imagen 5 sobre el canal verde permite lograr tonalidades más naturales.}
        % \end{figure}